\documentclass[11pt]{article}

\usepackage[margin=1.05in]{geometry}
\usepackage{amsmath,amssymb,amsthm,mathtools}
\usepackage{booktabs,array}
\usepackage{microtype}
\usepackage[hidelinks]{hyperref}

\allowdisplaybreaks
\setlength{\parindent}{1.2em}
\setlength{\parskip}{0.15em}

\newtheorem{theorem}{Theorem}[section]
\newtheorem{lemma}[theorem]{Lemma}
\newtheorem{proposition}[theorem]{Proposition}
\newtheorem{corollary}[theorem]{Corollary}
\theoremstyle{definition}
\newtheorem{remark}[theorem]{Remark}

\newcommand{\per}{\operatorname{per}}
\newcommand{\pcap}{\operatorname{Cap}}
\newcommand{\e}{\mathrm e}
\newcommand{\Rge}{\mathbb R_{\ge0}}

\title{A Codimension-One Entropy Refinement for Dittert's Conjecture in Dimension Seven}
\author{Extension note to the dimensions $8$--$15$ manuscript}
\date{24 July 2026}

\begin{document}
\maketitle

\begin{abstract}
This note gives a proposed extension of the Hall-deficit proof of Dittert's conjecture from dimensions $8$--$15$ to dimension $7$.  The new point is a geometric-mean lower bound for the cofactors of a selected column of a doubly stochastic matrix.  Applied twice, it strengthens the sharp permanent bound for a zero rectangle of codimension one: besides the sharp constant, the bound retains the entropy of both marginal distributions of the critical core.  The only two-sided Hall configurations not covered by the existing scalar method in dimension $7$ have codimension one.  Critical-core stationarity and a conditional product-energy estimate then reduce those configurations to three exact bivariate polynomial inequalities.  Their positivity on a rational square is certified by exact Bernstein subdivision.  The remaining Hall configurations are covered by the original second-order and one-sided arguments with sharpened constants.  No floating-point assertion is used.
\end{abstract}

\section{Statement and inherited notation}

For a nonnegative $n\times n$ matrix $A$ whose entries sum to $n$, write $r_i$ and $c_j$ for its row and column sums and put
\[
 \Phi(A)=\prod_i r_i+\prod_jc_j-\per(A),
 \qquad
 \gamma_n=\frac{n!}{n^n}.
\]
We use the arbitrary-dimensional parts of the accompanying manuscript: positive global maximizers are uniform; a hypothetical boundary maximizer has deficits
\[
 P:=\prod_i r_i=1-\rho,
 \qquad Q:=\prod_jc_j=1-\sigma,
 \qquad H:=\per(A)=\gamma_n-\delta,
\]
with
\[
 0\le \rho,\sigma,
 \qquad \rho+\sigma\le\delta\le\gamma_n;
 \tag{1.1}\label{eq:budget}
\]
and a critical Hall pair produces a dilation $S=A/(1-\tau)$ dominating a doubly stochastic matrix $B$ with a forced zero rectangle.

We also use
\[
 v_k=\frac{k!}{k^k}\quad(k\ge1),\qquad v_0=1,
\]
and the sharp zero-rectangle bound
\[
 \per(B)\ge L^{\sharp}_{n,a,b}
 =\frac{v_{n-a}v_{n-b}}{v_{n-a-b}}.
 \tag{1.2}\label{eq:sharp}
\]

The purpose of the note is to prove the following extension.

\begin{theorem}\label{thm:n7}
For every nonnegative $7\times7$ matrix $A$ whose entries sum to $7$,
\[
 \Phi(A)\le 2-\frac{7!}{7^7}.
\]
Equality holds only at $A=J_7/7$.
\end{theorem}

The interior classification is already dimension-free.  It therefore suffices to exclude a boundary global maximizer when $n=7$.

\section{The new permanent inequality}

The missing information in the sharp bound \eqref{eq:sharp} is the nonuniformity of the mass that crosses the critical core.  The following cofactor inequality recovers it.

\begin{lemma}[Geometric cofactor entropy]\label{lem:cofactor}
Let $M$ be an $m\times m$ doubly stochastic matrix, let
\[
 \xi=(\xi_1,\ldots,\xi_m)
\]
be one selected column, and let
\[
 d_i=\per M(i\mid \xi)
\]
be the permanent of the submatrix obtained by deleting row $i$ and the selected column.  Then
\[
 \prod_{i=1}^m d_i^{\xi_i}
 \ge
 v_m\exp\!\left\{
 \frac12\sum_{i=1}^m\left(\xi_i-\frac1m\right)^2
 \right\}.
 \tag{2.1}\label{eq:cofactor}
\]
The usual convention is used for zero weights, and the boundary cases follow by continuity.
\end{lemma}

\begin{proof}
Write $Y$ for the other $m-1$ columns and, initially assuming $0\le\xi_i<1$, normalize its rows by
\[
 n_{ij}=\frac{y_{ij}}{1-\xi_i}.
\]
Let
\[
 L_i(x)=\sum_{j=1}^{m-1}n_{ij}x_j,
 \qquad
 p_i=\per N_{-i}.
\]
Then
\[
 d_i=\left(\prod_{k\ne i}(1-\xi_k)\right)p_i.
 \tag{2.2}\label{eq:di-pi}
\]

For positive parameters $t_i$, consider
\[
 q_t(x)=\sum_{i=1}^m\xi_it_i\prod_{k\ne i}L_k(x).
 \tag{2.3}\label{eq:qt}
\]
It is real stable, because it is obtained by differentiating in $z$ at $z=0$ the stable polynomial
\[
 \prod_{i=1}^m\bigl(L_i(x)+\xi_it_i z\bigr).
\]
Weighted AM--GM first gives
\[
 q_t(x)
 \ge
 \left(\prod_i t_i^{\xi_i}\right)
 \prod_k L_k(x)^{1-\xi_k}.
\]
A second weighted AM--GM inequality gives
\[
 L_k(x)\ge\prod_jx_j^{n_{kj}}.
\]
Since $(1-\xi_k)n_{kj}=y_{kj}$ and every column of $Y$ sums to one,
\[
 q_t(x)\ge
 \left(\prod_i t_i^{\xi_i}\right)x_1\cdots x_{m-1}.
\]
Thus
\[
 \pcap(q_t)\ge\prod_i t_i^{\xi_i}.
\]
Square-free coefficient extraction for stable polynomials yields
\[
 \sum_i\xi_i t_i p_i
 \ge v_{m-1}\prod_i t_i^{\xi_i}.
 \tag{2.4}\label{eq:weighted-pi}
\]
Putting $t_i=p_i^{-1}$, with the zero-weight terms omitted and then using continuity, gives
\[
 \prod_i p_i^{\xi_i}\ge v_{m-1}.
 \tag{2.5}\label{eq:pi-gm}
\]

Taking the $\xi$-weighted geometric mean in \eqref{eq:di-pi},
\[
 \prod_i d_i^{\xi_i}
 \ge
 v_{m-1}\prod_k(1-\xi_k)^{1-\xi_k}.
 \tag{2.6}\label{eq:cofactor-intermediate}
\]
For $h(t)=(1-t)\log(1-t)$ one has $h''(t)=1/(1-t)\ge1$.  Hence
\[
 \sum_kh(\xi_k)
 \ge mh(1/m)+\frac12\sum_k(\xi_k-1/m)^2.
\]
Finally,
\[
 v_{m-1}\exp\{mh(1/m)\}=v_{m-1}\left(1-\frac1m\right)^{m-1}=v_m,
\]
and \eqref{eq:cofactor} follows.
\end{proof}

The codimension-one case of the critical zero rectangle now has a two-sided entropy correction.

\begin{theorem}[Codimension-one block entropy]\label{thm:block-entropy}
Let $B$ be an $n\times n$ doubly stochastic matrix with a zero block
\[
 B[R,C]=0,
 \qquad |R|=a,
 \qquad |C|=b,
 \qquad a+b=n-1.
\]
Put $I=R^c$ and $J=C^c$.  The core $X=B[I,J]$ has total mass one.  Let $\xi$ and $\eta$ be its row and column marginals and put
\[
 s_I=\sum_{i\in I}\left(\xi_i-\frac1{b+1}\right)^2,
 \qquad
 s_J=\sum_{j\in J}\left(\eta_j-\frac1{a+1}\right)^2.
\]
Then
\[
 \per(B)\ge
 v_{b+1}v_{a+1}\exp\!\left(\frac{s_I+s_J}{2}\right).
 \tag{2.7}\label{eq:block-entropy}
\]
At $s_I=s_J=0$, the constant is exactly the sharp bound \eqref{eq:sharp} with $p=1$.
\end{theorem}

\begin{proof}
Write $Y=B[I,C]$ and $Z=B[R,J]$.  Every perfect matching uses exactly one core edge.  If
\[
 d_i=\per Y_{-i},
 \qquad e_j=\per Z_{-j},
\]
then
\[
 \per(B)=\sum_{i\in I,\,j\in J}x_{ij}d_i e_j.
\]
Weighted AM--GM, using the core entries $x_{ij}$ as weights, gives
\[
 \per(B)\ge
 \prod_{i\in I}d_i^{\xi_i}
 \prod_{j\in J}e_j^{\eta_j}.
 \tag{2.8}\label{eq:block-amgm}
\]
The square matrix obtained by adjoining $\xi$ to $Y$ is doubly stochastic of order $b+1$.  Lemma~\ref{lem:cofactor} therefore bounds the first product in \eqref{eq:block-amgm} by $v_{b+1}\e^{s_I/2}$.  Likewise, adjoining $\eta^T$ as the last row of $Z$ and transposing bounds the second product by $v_{a+1}\e^{s_J/2}$.
\end{proof}

\section{Dimension-seven constants and the routine Hall pairs}

For the rest of the proof let
\[
 n=7,
 \qquad \gamma=\gamma_7=\frac{720}{117649},
 \qquad \eta=\frac1{47}.
\]
Define $\alpha_0=0$ and, for $1\le k\le6$,
\[
 \alpha_k=
 \begin{cases}
 14\left(\dfrac{k}{7}+\eta\right)
       \left(1-\dfrac{k}{7}-\eta\right),&2k<7,\\[2mm]
 \dfrac{2k(7-k)}7,&2k\ge7.
 \end{cases}
 \tag{3.1}\label{eq:alpha7}
\]
The relative-entropy subset proof from the main manuscript applies verbatim, because
\[
 \eta^2-\frac{\gamma}{14(1-\gamma)}
 =\frac{23263}{1808073127}>0,
 \qquad
 \frac37+\eta<\frac12,
 \tag{3.2}\label{eq:localization7}
\]
and $\alpha_k\le7/2$ for every $k$.

Let $a=|I^c|$, $b=|J^c|$, and $p=7-a-b$.  For two-sided pairs with $p\ge2$, use
\[
 T=\frac5{24},
 \qquad \lambda=\binom72-\binom73T=\frac{329}{24}.
\]
The Hall estimate gives $\tau<T$, since
\[
 T^2-\frac{7\gamma}{1-\gamma}
 =\frac{20185}{67351104}>0.
\]
The same alternating-binomial argument gives
\[
 (1-t)^7\ge1-7t+\frac{329}{24}t^2
 \qquad(0\le t\le5/24).
\]
Substitution in the second-order scalar margin, with
\[
 L=L^{\sharp}_{7,a,b},
 \qquad c=\frac{\alpha_a+\alpha_b}{p^2},
\]
shows that all six pairs with $p\ge2$ have positive margin.  The smallest is the pair $(a,b)=(1,1)$:
\[
 \frac{217499955444827479495}
 {21713873952987016007066496}>0.
 \tag{3.3}\label{eq:pge2-min}
\]

It remains to discuss $p=1$.  Up to transposition, the pairs are
\[
 (a,b)=(1,5),(2,4),(3,3).
 \tag{3.4}\label{eq:p1-pairs}
\]
Write
\[
 e=\sum_{i\in I^c}r_i-a,
 \qquad
 f=\sum_{j\in J^c}c_j-b,
 \qquad
 x=s_A(I^c,J^c).
\]
Then
\[
 \tau=e+f-x.
 \tag{3.5}\label{eq:tau-ef}
\]
If $e\le0$ or $f\le0$, only the positive excess is needed in the Hall estimate.  With
\[
 T=\frac3{20},
 \qquad \lambda=\binom72-\binom73T=\frac{63}{4},
\]
the six resulting second-order margins are positive.  Their minimum is
\[
 \frac{165035926385}{3897251304526692}>0.
 \tag{3.6}\label{eq:signmixed-min}
\]
Thus the only new case is \eqref{eq:p1-pairs} with
\[
 e>0,
 \qquad f>0.
 \tag{3.7}\label{eq:positive-ef}
\]
The subset bounds imply separately
\[
 0<e,f<\frac3{20}.
 \tag{3.8}\label{eq:ef-box}
\]

\section{One-sided Hall cuts in dimension seven}

For completeness, the one-sided part of the original proof also extends to $n=7$ once its two coarse numerical replacements are removed.

Put
\[
 d_0^2=\frac{7\gamma}{2(1-\gamma)}.
\]
Then $d_0<3/20$, and every row and column sum is at least $17/20$.  Moreover
\[
 \frac{(1-\gamma)(17/20)^2}{2}>\frac13.
 \tag{4.1}\label{eq:energy-coeff}
\]
Repeating the row-product energy proof therefore gives
\[
 \rho>\frac13\sum_i\left(1-\frac1{r_i}\right)^2.
 \tag{4.2}\label{eq:row-energy7}
\]
The column analogue is identical.

In the small-entropy regime, the exact certificate replacing the coefficient $2/5$ is
\[
 \frac13(1-\gamma)^2\gamma-\frac{49}{2}\kappa_7^2
 =
 \frac{22176489486262543850191055}
 {20672701805255574480144334848}>0,
 \tag{4.3}\label{eq:small7}
\]
where $\kappa_7=v_6G(6)=15625/2519424$.

For the large-entropy regime, retain the exact dependence on $\tau$ instead of replacing $(1-\tau)^{7/2}$ by a fixed decimal.  It is enough to prove, for $0\le t\le3/20$,
\[
 \frac23(1-\gamma)^2(1-t)^6
 -\gamma(7-21t+35t^2)^2>0.
 \tag{4.4}\label{eq:large7-poly}
\]
Every Bernstein coefficient of the polynomial in \eqref{eq:large7-poly} on $[0,3/20]$ is positive; the minimum is
\[
 \frac{155130361249919329}
 {1328763571296000000}>0.
 \tag{4.5}\label{eq:large7}
\]
Indeed, the exact one-sided inequality before squaring is
\[
 (1-\gamma)\sqrt{\frac{2\gamma}{3}}
 (1-\tau)^{5/2}
 >
 \gamma\frac{1-(1-\tau)^7}{\tau}.
\]
On this interval,
\[
 (1-\tau)^{5/2}\ge(1-\tau)^3,
 \qquad
 \frac{1-(1-\tau)^7}{\tau}
 \le7-21\tau+35\tau^2,
\]
which reduces the claim to \eqref{eq:large7-poly}.  Consequently all one-sided Hall cuts are excluded.

\section{The hard codimension-one configurations}

We now assume \eqref{eq:p1-pairs}, \eqref{eq:positive-ef}, and \eqref{eq:ef-box}.  Put
\[
 R=I^c,
 \quad C=J^c,
 \quad m=|I|=7-a=b+1,
 \quad \ell=|J|=7-b=a+1,
 \quad d=\min\{a,b\}.
\]

\subsection{Core saturation and stationarity}

Because $p=1$, both $B[I,J]$ and $S[I,J]$ have total mass one.  Since $B\le S$, one has entrywise saturation
\[
 B[I,J]=S[I,J]=\frac{A[I,J]}{1-\tau}.
 \tag{5.1}\label{eq:core-saturation}
\]
Let $\xi$ and $\eta$ be the row and column marginals of this core and put
\[
 s=s_I+s_J.
 \tag{5.2}\label{eq:s-total}
\]

For a permutation, the number of selected entries in $I\times J$ exceeds the number selected in $R\times C$ by exactly one.  Let $H_k$ be the total permanent weight of permutations using exactly $k$ entries of $R\times C$, so that
\[
 H=\sum_{k=0}^dH_k.
\]
Scale the critical core $A[I,J]$ by a parameter and globally renormalize to keep total mass seven.  Stationarity at the maximizer gives the exact identity
\[
 P\left(\sum_{i\in I}\frac{\xi_i}{r_i}-1\right)
 +Q\left(\sum_{j\in J}\frac{\eta_j}{c_j}-1\right)
 =\frac{\tau H+\sum_{k=1}^d kH_k}{1-\tau}.
 \tag{5.3}\label{eq:core-stationarity}
\]
Indeed, differentiating the permanent contributes $\sum_k(k+1)H_k$, while global normalization contributes $-(1-\tau)\Phi(A)$.

Let $H_0$ now denote the contribution of the permutations avoiding $R\times C$.  Theorem~\ref{thm:block-entropy}, \eqref{eq:core-saturation}, and $\tau\le e+f$ give
\[
 H_0\ge L(1-\tau)^7\e^{s/2}
 \ge L(1-e-f)^7\e^{s/2},
 \qquad
 L=v_{b+1}v_{a+1}.
 \tag{5.4}\label{eq:H0-lower}
\]
Since $\sum_{k\ge1}kH_k\le d(H-H_0)$, the right side of \eqref{eq:core-stationarity} is at most
\[
 \frac{(e+f+d)H-dH_0}{1-e-f}.
 \tag{5.5}\label{eq:stationarity-upper}
\]

\subsection{Conditional product energy}

Define the groupwise AM--GM envelopes
\[
 P_a(e)=\left(1+\frac ea\right)^a
        \left(1-\frac e{m}\right)^m,
 \qquad
 P_b(f)=\left(1+\frac fb\right)^b
        \left(1-\frac f{\ell}\right)^\ell.
 \tag{5.6}\label{eq:PaPb}
\]
Then $P\le P_a(e)$ and $Q\le P_b(f)$.  Put
\[
 u_r=P_a(e)-P,
 \qquad
 u_c=P_b(f)-Q.
\]

\begin{lemma}[Conditional inverse-energy]\label{lem:conditional-energy}
On the box \eqref{eq:ef-box},
\[
 \sum_{i\in I}\left(
 \frac1{r_i}-\frac1m\sum_{h\in I}\frac1{r_h}
 \right)^2\le3u_r,
 \tag{5.7}\label{eq:conditional-row}
\]
and the analogous column inequality is bounded by $3u_c$.
\end{lemma}

\begin{proof}
Every row and column sum is at least $17/20$.  Set $q=1-e/m$, so $q\ge17/20$.  For $r,q\ge17/20$ one has
\[
 \frac rq-1-\log\frac rq
 \ge\frac{(17/20)^2}{2}
       \left(\frac1r-\frac1q\right)^2.
 \tag{5.8}\label{eq:group-scalar}
\]
This follows by putting $y=1-q/r$ and writing the left side as
\[
 \int_0^y\frac{t}{(1-t)^2}\,dt.
\]
The appropriate groupwise linear terms cancel.  Hence
\[
 \log\frac{P_a(e)}P
 \ge\frac{(17/20)^2}{2}
 \sum_{i\in I}\left(\frac1{r_i}-\frac1q\right)^2
 \ge\frac{(17/20)^2}{2}V_I,
\]
where $V_I$ denotes the left side of \eqref{eq:conditional-row}.  Since
\[
 \begin{aligned}
 P_a-P
 &\ge P\log\frac{P_a}{P}\\
 &\ge(1-\gamma)\log\frac{P_a}{P}\\
 &>\frac13V_I
 \end{aligned}
\]
by \eqref{eq:energy-coeff}, the result follows.  The column proof is the same.
\end{proof}

\subsection{Reduction to a scalar inequality}

Put
\[
 t=e+f,
 \qquad
 D=2-P_a(e)-P_b(f),
 \qquad
 K=\gamma-D,
 \qquad
 u=\delta-D.
 \tag{5.9}\label{eq:K-u}
\]
The joint budget \eqref{eq:budget} gives
\[
 u\ge u_r+u_c\ge0,
 \qquad H=K-u.
 \tag{5.10}\label{eq:u-budget}
\]
Also put
\[
 A_0=L(1-t)^7,
 \qquad
 E=\frac e{m-e},
 \qquad
 F=\frac f{\ell-f}.
 \tag{5.11}\label{eq:A0EF}
\]
By the harmonic-mean inequality and Lemma~\ref{lem:conditional-energy}, the left side of \eqref{eq:core-stationarity} is at least
\[
 P_a(e)E+P_b(f)F-u_rE-u_cF-\sqrt{3su}.
 \tag{5.12}\label{eq:stationarity-lower}
\]
On the other hand, \eqref{eq:stationarity-upper}, \eqref{eq:H0-lower}, and \eqref{eq:u-budget} give the upper bound
\[
 \frac{(t+d)K-dA_0}{1-t}
 -\frac{(t+d)u+dA_0(\e^{s/2}-1)}{1-t}.
 \tag{5.13}\label{eq:stationarity-upper2}
\]
Because
\[
 E,F\le\frac{t+d}{1-t}
 \qquad\text{and}\qquad u_r+u_c\le u,
\]
the linear product-loss terms in \eqref{eq:stationarity-lower} are absorbed by the negative $u$ term in \eqref{eq:stationarity-upper2}.  Thus, with
\[
 \mathcal B=
 P_a(e)E+P_b(f)F
 -\frac{(t+d)K-dA_0}{1-t},
 \tag{5.14}\label{eq:Bscalar}
\]
stationarity requires
\[
 \mathcal B\le\sqrt{3su}.
 \tag{5.15}\label{eq:Benergy}
\]
In particular, if $\mathcal B>0$, then
\[
 su\ge\frac{\mathcal B^2}{3}.
 \tag{5.16}\label{eq:su-lower}
\]

The permanent estimate \eqref{eq:H0-lower} simultaneously gives
\[
 u+A_0\e^{s/2}\le K.
 \tag{5.17}\label{eq:per-budget}
\]
If $A_0>K$, this is already impossible.  Otherwise, using $\e^{s/2}\ge1+s/2$ and AM--GM,
\[
 su\le\frac{(K-A_0)^2}{2A_0}.
 \tag{5.18}\label{eq:su-upper}
\]
Consequently, a contradiction follows whenever
\[
 \mathcal B>0
 \qquad\text{and}\qquad
 2A_0\mathcal B^2>3(K-A_0)^2.
 \tag{5.19}\label{eq:master-cert}
\]

\subsection{Exact Bernstein certificate}

All quantities in \eqref{eq:master-cert} are rational functions of $e,f$.  Set
\[
 Q_0=(m-e)(\ell-f)(1-e-f),
 \qquad
 \widetilde{\mathcal B}=Q_0\mathcal B,
\]
and
\[
 \mathcal C
 =2A_0\widetilde{\mathcal B}^{\,2}
 -3(K-A_0)^2Q_0^2.
 \tag{5.20}\label{eq:Cpoly}
\]
The functions $A_0-K$, $\widetilde{\mathcal B}$, and $\mathcal C$ are bivariate polynomials with rational coefficients.

Exact Bernstein subdivision on
\[
 0\le e,f\le\frac3{20}
\]
proves, for each pair in \eqref{eq:p1-pairs}, that every point satisfies one of the following alternatives:
\[
 K<0,
 \qquad A_0-K>0,
 \qquad
 \widetilde{\mathcal B}>0\ \text{ and }\ \mathcal C>0.
 \tag{5.21}\label{eq:bernstein-alternatives}
\]
The first alternative contradicts $H=K-u\ge0$; the second contradicts \eqref{eq:per-budget}; and the third is \eqref{eq:master-cert}.  The exact subdivision audit is
\[
\begin{array}{c|ccc}
 (a,b)&\text{tree nodes}&A_0-K\text{ leaves}&
 \widetilde{\mathcal B},\mathcal C\text{ leaves}\\ \hline
 (1,5)&39&10&10\\
 (2,4)&35&9&9\\
 (3,3)&47&11&13
\end{array}
\tag{5.22}\label{eq:bernstein-table}
\]
A positive Bernstein coefficient lower-bounds the polynomial on its box, and every subdivision is the exact rational de Casteljau transformation.  Thus \eqref{eq:bernstein-alternatives} is a continuum certificate, not a mesh computation.

This excludes the last three critical Hall configurations.

\section{Completion}

A hypothetical boundary maximizer has $\tau>0$.  Two-sided pairs with $p\ge2$ are excluded by \eqref{eq:pge2-min}; codimension-one pairs with a nonpositive excess are excluded by \eqref{eq:signmixed-min}; the remaining codimension-one pairs are excluded by Sections~5.1--5.4; and one-sided pairs are excluded by Section~4.  Hence no boundary global maximizer exists in dimension seven.  The dimension-free interior classification then implies Theorem~\ref{thm:n7}.

\begin{remark}
The conceptual gain is not merely a better numerical constant.  The sharp block bound records only the sizes $(a,b,p)$.  At $p=1$, the critical core is itself a probability coupling, and its two marginal deviations are precisely the variables that couple to the stationarity equations.  Theorem~\ref{thm:block-entropy} keeps those variables rather than discarding them.  That is why the formerly negative scalar margins become coercive after stationarity is imposed.
\end{remark}

\end{document}
